\documentclass[11pt,a4paper]{article}
\usepackage[margin=24mm]{geometry}
\usepackage{amsmath,amssymb,amsthm,fontspec,unicode-math,booktabs,longtable,array,hyperref}
\setmainfont{STIXGeneral}\setmathfont{STIX Two Math}\setmonofont{Menlo}
\hypersetup{hidelinks}\setlength{\emergencystretch}{2em}
\newtheorem{theorem}{Theorem}\newtheorem{lemma}[theorem]{Lemma}
\newtheorem{proposition}[theorem]{Proposition}\newtheorem{corollary}[theorem]{Corollary}
\newcommand{\gen}[1]{\langle #1\rangle}\newcommand{\spn}[1]{\langle #1\rangle}
\newcommand{\Z}{\mathbb Z}\newcommand{\F}{\mathbb F}\newcommand{\one}{\mathbf1}\newcommand{\Aut}{\operatorname{Aut}}
\providecommand{\tightlist}{\setlength{\itemsep}{0pt}\setlength{\parskip}{0pt}}

\title{Regular halving, cubic toroids and the limits of preservation}
\author{Human Compression Project}\date{}
\begin{document}\maketitle
\begin{abstract}
We give a corrected account of the principal results in the inspected
preprint on generalized halving by Piedade and Tranchida. The literal
coset identification holds without residual connectivity of the output.
The latter property instead has an exact parity-shadow criterion, and a
finite affine regular hypertope shows that this criterion cannot be
omitted. For cubic toroids we determine the actual period lattices of
the halving groups, compare all displayed presentations, and prove every
allowable geometric branch and diagram. The source's constructions and
retained classical inputs are distinguished from the corrective proofs.
\end{abstract}

\section{The complete corrected contract}\label{A.contract}
The source throughout is arXiv:2405.19050v1, dated 29 May 2024~\cite{PT}.
Its 43 named items and the unnumbered classification in Figures 5--8
are routed individually in the accompanying source map. Repeated
introductory statements are not additional results. The complete
Topic C common paper~\cite{OldCommon} and its supplement~\cite{Common}
are both dependencies of this application. The former is retained
in full, including its binary-module arguments; its cost cannot be
replaced by only the particular theorems invoked here.

The following distinctions determine the mathematical contract.
Construction 3.6 is always a geometry and preserves thinness and the
natural chamber-transitive action under its stated hypotheses.
Residual connectivity requires (PC), and full residue transport then
holds. Thus the unrestricted residual-connectivity clauses in Theorem A,
Proposition 3.13 and Theorem 3.18 are false. The complete regular
counterexample below proves that regular thin input does not remove
this obstruction. Theorem B and the literal assertion of Theorem 4.4,
which identify the geometric construction with its halving-group coset
system, remain true by the maximal-stabilizer proof in the supplement.
A coset identification alone does not assert marked subgroup IP.
The introductory flag-transitivity converse in Theorem A is also false:
the final section gives a connected nonbipartite rank-three example whose
halving is flag-transitive although its original geometry is not. This
example satisfies (PC), so it is a separate obstruction, and it does not
affect the forward action statement in Proposition 3.14.

The general criterion for retaining another leaf is corrected from
``both endpoint shadows'' to ``at least one endpoint shadow'' over a
bipartite edge shadow. The source criterion is valid for a
flag-transitive input, since the endpoints then have equal shadow
bipartiteness. The supplement gives the complete five-cycle example
separating these assertions. Source Lemma 4.6 has an impossible literal
antecedent in the regular setting, rather than a counterexample to its
logical implication; the exact induced-neighbor coset criterion replaces
it. Source Examples 4.8 and 4.9 are proved there with their hypotheses
kept distinct.

For cubic toroids the classical input is the faithful cubical
realization and classification retained in source Theorem 5.1 from
McMullen--Schulte~\cite{MS}. The affine calculation below proves exact
marked groups and complete corrected presentations for every listed
parameter. Certain single extra relations in Propositions 5.8 and
5.11--5.13 present groups larger than intended. The weaker assertion
that the intended group is a quotient satisfying those relations
remains true. The final geometric proof supplies every nondegeneracy,
parity-compatibility and exhaustion argument for Figures 5--8,
independently of those insufficient presentations.

The borrowed bipartite theorem, source Theorem 4.1, is used only through
its geometric, faithful-action and regular conclusions proved in the
common papers. Its arbitrary nonthin full-automorphism semidirect-product
assertion is a retained borrowed boundary, not a newly recovered theorem
or a premise of these arguments. In the regular case the original
endpoint involution $a$ supplies the complement to the index-two color
subgroup, so the splitting used here is explicit. The informal remark
that the hemicube phenomenon seems sporadic is left as an observation;
no classification or open-question solution is claimed. The later
journal body is outside the inspected full-text boundary.

The arguments have the following provenance. The affine counterexample
is the accepted phase30 proof, reproduced completely below with its
finite certificate. The period-intersection and geometric-branch
proofs are the separately reviewed N014 packets, reproduced here with
only section and label changes. The common supplement similarly reuses
the phase27, stabilizer and further-leaf proofs. This assembly supplies
an organized dependency and source contract, not a new claim of priority
or a new authorship claim for inherited arguments.

\section{A finite affine construction and its
hypotheses}\label{a-finite-affine-construction-and-its-hypotheses}

Let M be a finite, thin, residually connected rank-three incidence
geometry with types called vertices, edges and faces. Let K act
regularly on its chambers. Assume the following concrete map conditions:
every edge has two distinct incident vertices and two distinct incident
faces; every incident vertex-face pair has exactly two common edges; and
an edge common to a vertex and a face makes that vertex-face pair
incident. These conditions hold for the seed below and are checked
directly, rather than inferred merely from subgroup IP.

Write \(\Omega\) for the edge set. For a seed vertex v let \(B_v\) be
its incident-edge set, and for a seed face f let \(D_f\) be its
incident-edge set. Let C be a K-invariant binary linear code in
\(\F_2^\Omega\), and put \(N=\F_2^\Omega/C\). Write \(q_e\) for the
image of the coordinate vector at e. For any block B among the \(B_v\)
and \(D_f\) put \(W_B=\operatorname{span}\{q_e:e\in B\}\). Require:

\begin{enumerate}
\def\labelenumi{\arabic{enumi}.}
\tightlist
\item
  All \(q_e\) are nonzero and distinct.
\item
  For every block B, \(q_e\in W_B\) if and only if \(e\in B\).
\item
  For every incident seed vertex-face pair (v,f),
  \(W_{B_v}\cap W_{D_f}=\operatorname{span}\{q_e:e\in B_v\cap D_f\}\).
\end{enumerate}

The quotient columns span N by construction. No assumption about
intersections for nonincident seed vertex-face pairs is needed.

Define a typed incidence system \(\Gamma\) as follows. Its type-0
objects are vectors \(n\in N\). Its type-1 objects are pairs
\((e,n+\langle q_e\rangle)\), identified with the graph edge
\(\{n,n+q_e\}\). Its type-2 objects are pairs \((v,n+W_{B_v})\); its
type-3 objects are pairs \((f,n+W_{D_f})\). The seed labels v and f
remain part of each object, even if two affine subspaces happen to
coincide.

Vertex-edge incidence is endpoint containment. A vertex is incident with
a higher object exactly when it belongs to its affine subspace. An edge
(e,L) is incident with (v,S), respectively (f,T), exactly when
\(e\in B_v\) and \(L\subseteq S\), respectively \(e\in D_f\) and
\(L\subseteq T\). Finally, (v,S) and (f,T) are incident exactly when
\textbf{v and f are incident in the seed and \(S\cap T\ne\varnothing\)}.
Distinct objects of the same type are not incident.

\subsection{Affine construction
theorem}\label{affine-construction-theorem}

Under these hypotheses \(\Gamma\) is a finite regular hypertope, with
regular group \(G=N\rtimes K\). It satisfies (B1) and (B2) at types
\{0,1\}.

Here regular hypertope means a thin, residually connected geometry with
a type-preserving group acting regularly on chambers. In fact G is its
full type-preserving automorphism group.

\textbf{Geometry and chambers.} A chamber is uniquely specified by a
vector n and a seed chamber (v,e,f): its four objects are n,
\((e,n+\langle q_e\rangle)\), \((v,n+W_{B_v})\), and \((f,n+W_{D_f})\).
Every flag extends to such a chamber. For a flag containing a vertex,
use its vector and extend the prescribed seed flag in M. For a flag
containing an edge but no vertex, choose either endpoint; all prescribed
higher affine subspaces contain this endpoint, and its prescribed seed
data form a flag. A flag containing only an incident higher pair has a
common vector by definition and two possible common seed edges. A single
higher object contains a vector and has an incident seed chamber. The
empty flag also extends. This proves the geometry axiom, with no
coset-Helly assumption.

\textbf{Thinness.} A panel missing type 0 has the two endpoints of its
edge as completions. A panel missing type 1 has its fixed vector and
incident seed vertex-face pair, hence exactly their two common edge
labels. They give distinct edges by hypothesis 1. A panel missing type 2
has the two possible seed vertices at its prescribed edge-face pair;
each determines exactly one affine coset through the fixed vector. The
dual argument treats type 3. Thus every panel has two completions.

\textbf{Residual connectivity.} The residue at a vector n is exactly M,
with types reordered as edges, vertices, faces: each seed higher label
has a unique affine coset through n.~Thus all flags containing a vertex
have the required connected residues. At an edge, the residue consists
of its two endpoints, its two seed vertices, and its two seed faces,
with every cross-type pair incident. It is the join of three two-point
sets; flags containing an edge therefore also have connected residues
whenever their corank is at least two.

The residue at a single higher object (v,S) contains the connected
Cayley graph on the affine space S with directions \(\{q_e:e\in B_v\}\).
Every remaining higher object in this residue contains a common vector
with S by the definition of incidence, and hence attaches to this
connected lower truncation. The same holds for a type-3 object. For an
incident pair (v,S),(f,T), its residue is precisely the vertex-edge
graph on the affine intersection \(S\cap T\) with directions
\(\{q_e:e\in B_v\cap D_f\}\). Hypothesis 3 makes it connected. Finally,
the global lower graph is connected because all \(q_e\) span N, and
every higher object contains lower vertices. These cases exhaust all
flags of corank at least two in rank four.

\textbf{Regularity.} Translation by N and the induced K-action on N and
on the seed labels preserve incidence. The chamber parametrization above
shows that \(N\rtimes K\) acts freely and transitively: translate the
vector, then use the regular action of K on seed chambers. A thin,
residually connected geometry of finite rank is chamber connected
(induct on rank inside residues, using connected incidence graphs). An
automorphism fixing a chamber fixes each adjacent chamber by thinness
and hence every chamber. Therefore any type-preserving automorphism is
the unique element of G with its prescribed action on a base chamber.

\textbf{The two leaf conditions.} Hypothesis 1 makes the lower graph
simple, which is (B1). An edge with endpoints in \(S=n+W_B\) has direction
\(q_e\in W_B\); hypothesis 2 says exactly that \(e\in B\), and then it
is incident with the higher object by definition. This is (B2).

\section{The seed and the binary
quotient}\label{the-seed-and-the-binary-quotient}

Use permutations of \{1,\ldots,7\}:

\[c=(1\,2)(3\,4)(6\,7),\quad b=(1\,5)(2\,3)(6\,7),\quad d=(1\,3)(2\,4)(6\,7).\]

Let \(K=\langle c,b,d\rangle\). The exact checker obtains
\(|K|=120\), \(|cb|=|bd|=5\),
\(cd=dc\), and all 64 subgroup-IP meets. The permutation group is
\(A_5\times C_2\), although that name is not needed in the proof. Take
the seed vertex, edge and face objects to be the left cosets of
\(\langle b,d\rangle\), \(\langle c,d\rangle\) and
\(\langle c,b\rangle\), with incidence defined by nonempty intersection.
The checker enumerates this incidence system itself: 12 vertices, 30
edges, 12 faces, 355 flags and 120 chambers; every flag extends to a
chamber, every panel has size two, every residue of corank at least two
is connected, and K is regular on chambers. Thus this is an actual
regular map of type \{5,5\}.

Each vertex-star and each face-boundary contains five edges. There are
60 incident vertex-face pairs, each with exactly two common edges;
nonincident pairs have no common edge. Define C to be the span of the
characteristic vectors of these 24 five-edge blocks in \(\F_2^{30}\).
This definition is K-invariant because K permutes the two block
families; the checker also verifies invariance under c,b,d explicitly.

Exact binary row reduction gives \(\dim C=17\), \(\dim N=13\), and the following
complete small certificate:

\begin{itemize}
\tightlist
\item
  30 nonzero, pairwise distinct quotient columns;
\item
  24 block spans of dimension four;
\item
  all 720 column-in-block-span membership tests satisfy hypothesis 2;
\item
  all 60 incident block-span intersections equal their two-column spans,
  of dimension two.
\end{itemize}

The concrete permutation marking, coset labels, code basis, quotient
columns and block lists are included in the checker output. These are
finite mathematical data. Neither the code computation nor its
certificate substitutes for the global geometry proof in the affine
construction theorem above.

Each five-column block has sum zero and rank four. Its only nonempty
linear dependence is therefore the sum of all five columns. In
particular its affine Cayley graph has a simple five-cycle and is
nonbipartite. At an incident higher pair, the common shadow has two
independent directions and hence is a square.

\section{Distinguished involutions and
sizes}\label{distinguished-involutions-and-sizes}

Choose the seed base chamber at the identity cosets, and let
\(e_0=\langle c,d\rangle\). Let a be translation by \(q_{e_0}\), and
regard b,c,d as elements of the complement K in \(G=N\rtimes K\). The
chamber-adjacent involutions of \(\Gamma\), in type order, are
(a,b,c,d). Indeed a exchanges the endpoints of its affine edge, b
exchanges the two seed edges at the prescribed vertex-face pair, c
changes the seed vertex, and d changes the seed face.

The edge stabilizer \(\langle c,d\rangle\) fixes \(q_{e_0}\), so a
commutes with c,d.~The two columns \(q_{e_0}\),\(q_{be_0}\) are
distinct, and \((ab)^2\) is their nonzero sum, a translation of order
two. Thus

\[|ab|=4,\quad |bc|=|bd|=5,\quad |ac|=|ad|=|cd|=2.\]

The diagram is a star with centre b and labels 4,5,5. In particular a is
a nondegenerate leaf of the required kind. Conjugates of a by K give all
column translations, so these four involutions generate G. This marking
has actual subgroup IP because it is the distinguished marking of the
regular hypertope just proved, not because an arbitrary involutory
marking is assumed to define a geometry.

The object and chamber counts are

\[|N|=8192,\qquad |G|=983040,\]

\[\text{object counts}=(8192,122880,6144,6144),\qquad\text{chambers}=983040.\]

The standard maximal parabolics have orders 120,8,160,160. For example a
higher-object stabilizer is its 16-element translation space extended by
its 10-element seed stabilizer. The base edge stabilizer has order
eight. An incident pair of higher objects has a square residue, with
eight chambers; hence there are 983040/8=122880 such pairs.

\section{The geometric failure after
halving}\label{the-geometric-failure-after-halving}

The global lower graph is nonbipartite, since each higher shadow already
contains a five-cycle. All individual higher shadows are nonbipartite,
so Construction 3.6 of the inspected source does not split any higher
object by parity. Its two new lower types are the two labelled copies of
N, with incidence \((n,0)\sim(m,1)\) precisely when n,m were adjacent in the
original lower graph. Incidence with each unsplit higher object is
affine containment in either copy. Higher-higher incidence is unchanged
here.

Fix any incident pair of higher objects. Its original common shadow is a
square. Its new rank-two residue is the canonical bipartite double cover
of that square, namely two disjoint squares. Thus the new incidence
geometry is not residually connected. There are exactly 122880 such
failed residues. This is a literal local description, not an inference
from a numerical graph test. The original affine construction has
already established geometry, thinness, residual connectivity,
regularity, (B1), (B2), and the leaf diagram independently.

For clarity, the transformed incidence system is still a geometry and is
thin: a flag of higher objects has a common original vector, and, for a
higher pair, a common adjacent original pair; each transformed flag
extends using the original chamber description. The diagonal action of G
on both labelled lower copies is regular on transformed chambers, again
indexed by oriented original edges and their two seed higher objects.
The rank-two residues above have more than one chamber-connected
component; a free and transitive action on the global chambers does not
repair those disconnected residues.

\section{A direct subgroup-IP failure of the halved
marking}\label{a-direct-subgroup-ip-failure-of-the-halved-marking}

Put \(a^{\prime}=aba\) and \(H=\langle a^{\prime},b,c,d\rangle\). In the
local group \(U_c=\langle a,b,c\rangle\), the subgroup
\(\langle b,c\rangle\) is dihedral of order ten and is transitive on the
five edges of the base face. The translation \(a^{\prime}b\) is
\(q_{e_0}+q_{be_0}\). Consider the orbit of the unordered pair
\(\{e_0,be_0\}\) under this transitive degree-five action. Its
graph is connected: all components have the same size dividing five, and
the existence of an edge rules out singleton components. Pair
differences along a connected graph span all even sums of its coordinate
vectors. Thus the conjugates of \(a^{\prime}b\) span the even sums of
the five face columns. Since the sum of all five columns is zero, these
even sums already span the whole four-dimensional space. Consequently
\(a\in\langle a^{\prime},b,c\rangle\) and

\[\langle a^{\prime},b,c\rangle=\langle a,b,c\rangle,\qquad |\langle a,b,c\rangle|=160.\]

The same argument with \(\langle b,d\rangle\) and the base vertex-star
gives

\[\langle a^{\prime},b,d\rangle=\langle a,b,d\rangle,\qquad |\langle a,b,d\rangle|=160.\]

In particular \(H=G\). Actual IP for the original regular marking gives

\[\langle a^{\prime},b,c\rangle\cap\langle a^{\prime},b,d\rangle=\langle a,b\rangle,\qquad |\langle a,b\rangle|=8.\]

But \(a^{\prime}\) and b are distinct commuting involutions, since
\((ab)^4=1\), and \(\langle a^{\prime},b\rangle\) has order four. It
does not contain a: in the dihedral group \(\langle a,b\rangle\) of
order eight, the subgroup \(\langle aba,b\rangle\) consists of the
elements having an even number of occurrences of a. Thus a is an
explicit witness to the failed meet. The final checker also constructs
both local semidirect groups and this strict intersection directly.

This obstruction is precisely the nonprincipal recovery family from the
general halving criterion: a is recovered after adjoining c or d
separately, but not before either is adjoined. Here the original
incidence system is an actual regular hypertope, unlike the earlier \(S_6\)
pregeometry separator.

\section{Source and dependency boundary of the regular
example}\label{source-and-dependency-boundary-of-the-regular-example}

The inspected primary is Piedade--Tranchida, arXiv:2405.19050v1 (29 May
2024). Its exact input identity is recorded with this assembly. The
source's broad Theorem A and its residual-connectivity assertion in
Proposition 3.13 / Theorem 3.18 fail on this example, now even with
regular thin input. In §4, the sentence immediately before Theorem 4.4
says the resulting object must be a regular hypertope; this claimed
consequence also fails, as does the assertion that the halved
distinguished generators must form a C-group.

The literal sentences of Theorem B and Theorem 4.4 assert a
coset-isomorphism, a narrower statement. This note does \textbf{not}
claim to refute that isomorphism. Indeed equality of the local new and
old higher parabolics is compatible with the transformed geometry still
having its asserted coset description while failing residual
connectivity. We do not conflate those assertions.

The revised journal article, Combinatorica 45 (2025), article 10, DOI
10.1007/s00493-024-00134-y, has a later revision date. Its full revised
body was not available in the bounded source audit. No claim about that
body, absence of a subsequent correction, or novelty is made here.

The mathematical construction uses elementary permutation-group closure,
binary linear algebra, the explicitly proved affine construction
theorem, and the elementary chamber-connectivity property of finite-rank
thin residually connected geometries. It uses no subgroup catalogue, no
presumed quotient-IP theorem, and no full enumeration of the
983040-element group or the global incidence system. The retained parity
terminology and source identification are separate from the direct
square-residue and local-subgroup proofs above.

\section{The complete finite
certificate}\label{the-complete-finite-certificate}

The complete program \texttt{phase30\_affine\_map\_check.py}
produces the saved certificate. It uses Python's standard
library only. The whole program and its saved exact output accompany
this paper. The replay command is recorded in the assembly guide.

The completed check enumerates all 355 seed flags, their extensions and
residues, all 64 seed subgroup meets, all 720 column memberships, all 60
literal four-element span intersections, 72 code-generator and 51
code-basis invariance tests, and the two old and two changed local
semidirect groups of order 160. It verifies all six marked product
orders and the explicit excess-intersection witness a. It also
constructs the eight-vertex doubled square and its two four-vertex
components. The JSON includes both 30-coordinate quotient
representatives and compressed 13-bit columns, all block labels, a basis
of the code, and the induced generator action on the quotient.

The independent parent checker
\texttt{phase30\_regular\_halving\_root\_check.py} uses
\(A_5\times C_2\) product coordinates instead of seven-point
permutations and separately checks the same seed geometry, binary
hypotheses and local subgroup failure. Neither script enumerates the
whole 983040-element group or its global incidence system. Global
geometric conclusions and their derived counts depend on the affine
construction and local recovery proofs above. The earlier exploratory
probe is excluded from the final authored export manifest.

The translation lattices of cubic toroids give a direct way to compute
their halving groups. One first takes the subgroup generated by the
new reflections in the infinite affine group, and then intersects its
translation lattice with the original period lattice. This gives all
the groups in Propositions~5.8--5.13 of the inspected preprint~\cite{PT}
by one calculation. It also detects missing relations in several of
the claimed presentations. We distinguish the calculation of a marked
group from the geometric conditions required to perform a halving.

\section{Affine notation and the two subgroups}
Let $n\ge3$. Let $B_n$ be the signed permutation group on $\Z^n$,
and let $B_n^+$ be its subgroup with an even number of negative signs.
Thus $|B_n|=2^nn!$ and $|B_n^+|=2^{n-1}n!$.
Let $s_i$ exchange coordinates $i,i+1$, let $z_i$ change the sign
of coordinate $i$, and let $\tau_v$ denote translation by $v$.
We multiply affine transformations by composition, with the rightmost
factor acting first. The usual cubic marking of
$W=\Z^n\rtimes B_n$ is
\[
 a=\tau_{e_1}z_1,\qquad s_1,\ldots,s_{n-1},\qquad z_n.
\]
We retain its standard affine Coxeter presentation
$[4,3^{n-2},4]$. This is the same classical realization used in
Theorem~5.1 of~\cite{PT}. Put
\[
 D=\{v\in\Z^n:\textstyle\sum_i v_i\equiv0\pmod2\},\quad
 x=as_1a=\tau_{e_1+e_2}z_1z_2s_1,\quad
 y=z_ns_{n-1}z_n=z_{n-1}z_ns_{n-1}.
\]
The new markings after one and two endpoint operations generate
\[
 K=\gen{x,s_1,\ldots,s_{n-1},z_n},\qquad
 J=\gen{x,s_1,\ldots,s_{n-1},y}.
\]

\begin{lemma}\label{NA.universal}
One has $K=D\rtimes B_n$ and $J=D\rtimes B_n^+$.
The Coxeter presentations of these groups are, respectively, the
diagram obtained by halving the first endpoint of
$[4,3^{n-2},4]$, and the diagram obtained by halving both endpoints.
In rank four the latter is a square of bonds labelled three;
in rank five it is the four-leaf diagram with all bonds labelled three.
\end{lemma}
\begin{proof}
The linear generators of $K$ generate $B_n$. Multiplying $x$ by
$z_1z_2s_1$ gives $\tau_{e_1+e_2}$. Its conjugates by $B_n$
are translations by all vectors $\pm e_i\pm e_j$, which span $D$.
Conversely every displayed generator belongs to $D\rtimes B_n$.
For $J$, the generators $s_i$ and $y$ generate $B_n^+$: the
adjacent swaps give all permutations, while $ys_{n-1}$ changes
two signs and its conjugates give every pair of sign changes.
Again $z_1z_2s_1\in B_n^+$, so $\tau_{e_1+e_2}\in J$.
The orbit of this vector under $B_n^+$ contains all
$\pm e_i\pm e_j$ when $n\ge3$. To change just one of its two
signs one changes the sign of a third, zero coordinate as well.
This proves the second equality.

Here is also a presentation argument, needed when comparing the
printed extra relations. In a Coxeter presentation with an endpoint
$a$ joined to $b$ by a bond four and commuting with the remaining
generators, replace $a$ by $x=aba$. The proposed new Coxeter
diagram duplicates the bonds of $b$ at $x$, and makes $x,b$ commute.
Its abstract Coxeter group has an involutory automorphism exchanging
$x,b$ and fixing the other generators. Adjoin an involution $a$
acting by this automorphism. Eliminate $x$ using $x=aba$.
The relation $(xb)^2=1$ becomes $(ab)^4=1$; the remaining
relations at $x$ are conjugates of those at $b$. All other
relations are exactly the old Coxeter relations. Thus the old
group is the semidirect product of the new Coxeter group by
$\gen a$, and the new group embeds as the kernel of the endpoint
parity character. Apply this argument first at $a$ in $W$, then
at $z_n$ in $K$. In the second application the character is the
parity of the number of negative signs. This gives precisely
the groups already identified and also covers $n=3,4$.
\end{proof}

\section{The period lattice is intersected, not merely squared}
For $s\ge2$, the three toroid period lattices are
\[
 \Lambda_1=s\Z^n,\qquad \Lambda_2=sD,\qquad
 \Lambda_n=s(2\Z^n+\Z\one),\qquad \one=(1,\ldots,1).
\]
These are the spans of the signed permutation orbits of
$se_1$, $s(e_1+e_2)$ and $s\one$, respectively, as in~\cite{PT}.
They are normal translation subgroups of $W$.
Let $\overline K$ and $\overline J$ denote the images of the two
marked groups in $W/\Lambda_k$.

\begin{theorem}\label{NA.lattice}
Set $L=D\cap\Lambda_k$. Then, with their displayed affine markings,
\[
 \overline K=(D/L)\rtimes B_n,\qquad
 \overline J=(D/L)\rtimes B_n^+.
\]
The exact lattice and its index in $D$ are as follows.
\[
\begin{array}{c|c|c}
\text{parameters}&L&[D:L]\\\hline
k=1,\ s\text{ even}&s\Z^n&s^n/2\\
k=1,\ s\text{ odd}&sD&s^n\\
k=2&sD&s^n\\
k=n,\ ns\text{ even}&s(2\Z^n+\Z\one)&2^{n-2}s^n\\
k=n,\ ns\text{ odd}&2s\Z^n&2^{n-1}s^n
\end{array}
\]
In particular $|\overline K|=2^nn![D:L]$ and
$|\overline J|=2^{n-1}n![D:L]$.
\end{theorem}
\begin{proof}
Restriction of the quotient homomorphism $W\to W/\Lambda_k$
to either group in Lemma~\ref{NA.universal} has kernel $D\cap\Lambda_k$.
The linear subgroup intersects this kernel trivially. It therefore
remains a complement, proving the two semidirect product statements.

A vector $sv$ has even coordinate sum automatically when $s$ is
even, and exactly when $v\in D$ when $s$ is odd. This gives
the first two rows. The third is immediate. Write an element of
$\Lambda_n$ as $s(2v+t\one)$. Its coordinate sum modulo two is
$nst$. If $ns$ is even there is no restriction; otherwise $t$
must be even, giving $2s\Z^n$.
Finally $[\Z^n:D]=2$ and
$[\Z^n:2\Z^n+\Z\one]=2^{n-1}$.
Multiplication of a full lattice by $s$ multiplies its index by
$s^n$. These facts give all the displayed indices.
\end{proof}

This is a complete group description, rather than only an order
calculation. It also supplies finite presentations uniformly.
Choose the basis of $D$
\[
 d_1=e_1+e_2,\quad d_2=e_2-e_1,\quad
 d_i=e_i-e_{i-1}\quad(3\le i\le n).
\]
Adjoin commuting translation generators $T_1,\ldots,T_n$ to a
standard presentation of the finite group $P=B_n$ or $B_n^+$.
For each generator $p$ of $P$ impose
\[
 pT_jp^{-1}=\prod_{i=1}^nT_i^{c_{ij}(p)},\qquad
 pd_j=\sum_i c_{ij}(p)d_i.
\]
For a basis of $L$ impose $\prod_iT_i^{\ell_i}=1$ whenever
$\sum_i\ell_i d_i$ is a member of that basis. Such a basis is
$se_i$ in the first row, $sd_i$ in the second and third,
$2se_1,\ldots,2se_{n-1},s\one$ in the fourth, and $2se_i$
in the fifth. All these vectors lie in $D$ under the stated
parity assumptions, so their coefficients are integers.
The conjugation relations move every word to the form $T^v p$.
The multiplication law and the evident affine realization then
show that these presentations give exactly $(D/L)\rtimes P$.
The marked element $x$ is $T_1z_1z_2s_1$, and the other marked
elements are the specified linear transformations. Thus the
presentation retains the actual marking as well as the abstract group.

\section{Comparison with the single extra relations}
The defect in a proposed extra relation can now be tested inside
the infinite groups, before taking any finite quotient. A pure
translation $\tau_v$ normally generates the translations whose
vectors lie in the integer span of the orbit of $v$ under the
linear complement. Translations conjugate one another trivially,
so there are no additional conjugates to supply missing periods.

For clarity we evaluate every family of extra words used in
Propositions~5.8--5.13. Put
\[
 u=as_1\cdots s_{n-1}z_n,
 \qquad u(X_1,\ldots,X_n)=(X_n+1,X_1,\ldots,X_{n-1}).
\]
Thus $u^n=\tau_\one$. The words for $k=n$, both before and
after the second halving, are $u^2$ written in the new generators.
This follows by moving $a$ past $s_i$ for $i\ge2$ in the first
case; in the second, move $z_n$ past $s_i$ for $i<n-1$ and
replace $z_ns_{n-1}z_n$ by $y$. The resulting word is
\[
 x s_2\cdots s_{n-1}s_1s_2\cdots s_{n-2}y.
\]
Its power $ns/2$, when integral, is $\tau_{s\one}$; its
power $ns$ in the odd case is $\tau_{2s\one}$.

For $k=1$, the single-halving word is
\[
 q=x s_2\cdots s_{n-1}z_ns_{n-1}\cdots s_2s_1
   =\tau_{e_1+e_2}z_2.
\]
The corresponding word after two halvings is
\[
 q'=x s_2\cdots s_{n-1}y s_{n-2}\cdots s_2s_1
    =\tau_{e_1+e_2}z_2z_n.
\]
Both square to $\tau_{2e_1}$. Thus their displayed powers
are $\tau_{se_1}$ for even $s$ and $\tau_{2se_1}$ for odd $s$.

For $k=2$, put
\[
 A=s_2\cdots s_{n-1}z_ns_{n-1}\cdots s_2=z_2,\qquad
 A'=s_2\cdots s_{n-1}y s_{n-2}\cdots s_2=z_2z_n.
\]
These concern the first and second operation, respectively.
The words $xAs_1A$ and $xA's_1A'$
are both $\tau_{e_1+e_2}$. The rank-four word printed separately
in Proposition~5.12 is $xs_2s_1ys_1s_2$; conjugating $y$
by $s_2s_1$ shows its linear part is $z_1z_2s_1$, so it has
the same value. Consequently the extra relator is the translation
$\tau_{s(e_1+e_2)}$ in every $k=2$ case.

\begin{proposition}\label{NA.comparison}
Interpret the listed extra relations in~\cite{PT} as complete
presentations, as their proofs do. They have the following exact
discrepancies. The factor is the order of the presented group
divided by the order of the intended marked group.
\[
\begin{array}{c|c|c}
\text{parameters}&\text{one halving}&\text{two halvings}\\\hline
k=1,\ s\text{ odd}&2^{n-1}&2^{n-1}\\
k=n,\ ns\text{ odd}&2^{n-1}&2^{n-1}\\
k=n,\ n\text{ even}&1&2\\
\text{all other cases}&1&1
\end{array}
\]
The table compares presentations of marked groups. Geometric
validity of the second operation is a separate hypothesis.
\end{proposition}
\begin{proof}
The signed permutation orbit of $e_1$ spans $\Z^n$ under
either complement, and that of $e_1+e_2$ spans $D$.
The orbit of $\one$ under $B_n$ spans
$E=2\Z^n+\Z\one$: changing one sign gives $2e_i$, and these
vectors together with $\one$ also generate every orbit vector.
Its orbit under $B_n^+$ instead spans
\[
 M=2D+\Z\one.
\]
Indeed differences from $\one$ are twice the characteristic
vectors of even subsets; the two-element subsets already span $D$.
If $n$ is odd, then $\one-e_i\in D$, so
$2e_i=2\one-2(\one-e_i)\in M$ and $M=E$.
If $n$ is even, $M$ has index two in $E$. To see this, every
generator $2e_i$ of $2\Z^n$ has the same coset modulo $2D$,
so the index is at most two. It is not one: an equality
$2e_1=2d+t\one$, with $d\in D$, forces $t=2h$ and then
$e_1=d+h\one$, contradicting parity of coordinate sums.

Let $R$ be the normal translation lattice generated by the
printed extra relation. For odd $s$ and $k=1$, one has
$R=2s\Z^n$, whereas $L=sD$. Their relative index is
$[sD:2s\Z^n]=2^{n-1}$. For odd $ns$ and $k=n$, one has
$R=2sE$ for either operation, whereas $L=2s\Z^n$;
again the relative index is $2^{n-1}$. For even $n$ and $k=n$,
the first operation gives $R=sE=L$, while the second gives
$R=sM$, of index two in $L=sE$.
The other rows give $R=L$ by the same orbit calculations.
Lemma~\ref{NA.universal} identifies the underlying Coxeter group
exactly, so its quotient by the displayed relation is
$(D/R)\rtimes P$. The factor is therefore exactly $[L:R]$,
not merely a lower bound from a finite quotient.
\end{proof}

For example, at $n=3,k=1,s=3$, the intended first halving group
has order $3^3\cdot48=1296$. The displayed presentation instead
has order $[D:6\Z^3]\cdot48=108\cdot48=5184$.
The accompanying small diagnostic constructs both affine groups
directly, verifies every Coxeter relation and the displayed extra
word, and obtains these orders. The all-parameter conclusions
above follow from the lattice proofs, not from this finite check.

\section{Attribution and the geometric interface}
The cubic tessellation, the three period lattices and the
halving operations are prior work. The parity characters,
the second isomorphism theorem, and the standard Coxeter
presentations are also established ingredients. Our contribution
here is their common application to the exact period intersection,
the complete presentations that it supplies, and the comparison
that identifies the missing relations in the inspected edition.
We make no claim of publication priority or about a later journal
revision. The assertions that particular iterated constructions
remain hypertopes, and the full list of their diagrams, are proved in the geometric sections that follow, with the nondegeneracy
and residue arguments explicitly supplied.

This note supplies the geometric arguments behind Section~5 and
Figures~5--8 of the inspected preprint of Piedade--Tranchida~\cite{PT}.
The constructions are theirs. We use the previously reviewed
component formulation of halving~\cite{Parity}, including its precise
residual-connectivity condition, and the separately reviewed affine
group calculation~\cite{Affine}. In particular a presentation or a
Coxeter diagram will not be used to infer residual connectivity.
All assertions about the source refer to arXiv:2405.19050v1,
not to an uninspected later journal edition.

\section{The cubical model and the retained halving theorem}
Let $n\ge3$, $s\ge2$, and $k\in\{1,2,n\}$. In the usual cubical
tessellation of $\mathbb R^n$ put
\[
 D=\{v\in\Z^n:\textstyle\sum v_i\equiv0\pmod2\},\qquad
 L=D\cap\Lambda_k,
\]
\[
 \Lambda_1=s\Z^n,\quad\Lambda_2=sD,\quad
 \Lambda_n=s(2\Z^n+\Z\one).
\]
A cubical face is a product of integral singleton coordinates and
unit intervals. Its dimension is its number of interval coordinates.
The original toroid has objects the $\Lambda_k$-orbits of faces of
dimensions $0,\ldots,n$, with incidence induced by containment.
We retain the classical assertion that this is the regular polytope
of Theorem~5.1 of~\cite{PT}, attributed there to
\cite[Theorem~6D4]{MS}. Its usual marking, acting on column
coordinates, is
\[
 a=\tau_{e_1}z_1,\quad s_1,\ldots,s_{n-1},\quad z_n,
\]
where $s_i$ swaps coordinates $i,i+1$ and $z_i$ changes the sign of
coordinate $i$. This is the standard affine realization, not a new
classification of cubic toroids.

We recall exactly the halving result used here. A selected vertex--edge
truncation is required to be a simple graph $Q$ (B1). For every other
object, incidence with an edge must be equivalent to incidence with
both endpoints (B2). For a flag $F$ of other objects let $Q_F$ be its
common graph shadow. If $Q$ is nonbipartite, the additional condition is
\[
 \tag{PC}
 \text{all singleton shadows }Q_x\ (x\in F)\text{ nonbipartite}
 \ \Longrightarrow\ Q_F\text{ nonbipartite}.
\]
In the component construction one takes a component of the bipartite
double cover of $Q$, and lifts every other object by the components of
the double cover of its shadow. The resulting incidence system is a
geometry; it is residually connected when $Q$ is bipartite, and, when
$Q$ is nonbipartite, exactly when PC holds. Thinness and chamber
transitivity are preserved. These statements, including the residue
isomorphisms used to compute the diagram, are the retained results of
\cite{Parity}. The regular coset identification and full type-preserving
automorphism group are supplied by~\cite{Coset}. Neither result permits
one to omit PC in a nonbipartite step.

Write $H_{ab}$ for the geometry after $a$ operations at the left end
and $b$ at the right end, with $H_{00}$ the original toroid. The notation
will only be used for the sequences proved to exist below. Choices of
the two equivalent leaves at a fork are identified by a type
permutation. A ``second operation at the left'' means the leaf
$(0,2)$ after the first operation $(0,1)$.

\section{Small translations and local cubical incidence}
The following elementary separation facts keep the quotient arguments
explicit. Put $R=\{\pm e_i\pm e_j:i\ne j\}$.

\begin{lemma}\label{NI.separation}
Every nonzero vector of $\Lambda_k$ or $L$ has a coordinate of absolute
value at least $s$. No two vertices of a unit cube are therefore
identified. Also:
\begin{enumerate}
\item $2e_i\in\Lambda_k$ or $2e_i\in L$ occurs exactly when
$(k,s)=(1,2)$.
\item If $r,r'\in R$ are distinct and $r-r'\in L$, then
$(k,s)=(1,2)$ or $(k,s)=(2,2)$. Conversely both parameter pairs
allow such a coincidence.
\item Suppose neither of these two parameter pairs occurs. If two
even vertices of one unit cube differ by a root modulo $L$, they
differ by that root literally. Their unique square is a face of
the cube. If two even vertices at root distance have a common odd
unit-step neighbor modulo $L$, that neighbor is one of the other
two corners of this square, after a common lift.
\end{enumerate}
The corresponding statements hold in the dual cubical grid of
cube centers, translated by $\frac12\one$.
\end{lemma}
\begin{proof}
All lattices are contained in $s\Z^n$. For (1), it remains only to
check $s=2$. The vector $e_i$ is neither in $D$ nor in
$2\Z^n+\Z\one$ for $n\ge3$.

For (2), every coordinate of $r-r'$ has absolute value at most two.
Thus $s=2$. If $r-r'$ is even coordinatewise, the supports of $r,r'$
are equal, and its nonzero support has size at most two. In $2D$ a
single coordinate $\pm2$ is excluded; a pair $\pm2,\pm2$ is allowed.
In $2(2\Z^n+\Z\one)$ neither a one-coordinate nor a two-coordinate
vector of this form occurs, since $n\ge3$. Passing to $L$ cannot
create a coincidence. Examples are $e_1+e_2$ and $-e_1+e_2$ for
$(1,2)$, and $e_1+e_2$ and $-e_1-e_2$ for $(2,2)$.

For (3), lift the cube and its two vertices $p,q$. In
$q-p-r\in L$, every coordinate has absolute value at most two.
For $s\ge3$ it is zero. For $s=2$, a nonzero even coordinate
can occur only on the two coordinates supporting $r$, and the
remaining allowed lattice $2(2\Z^n+\Z\one)$ contains no such
nonzero vector. Hence $q-p=r$. The square has exactly these two
varying coordinates, so is unique and contained in the cube.

For the common-neighbor assertion, first lift the neighbor $w$ and
one endpoint $p$ so that $w-p$ is a unit vector. The other incidence
gives $q-p\equiv\alpha+\beta\pmod L$, where $\alpha,\beta$ are
signed unit vectors. This sum is zero, a root, or $\pm2e_i$.
The zero case cannot equal a root modulo $L$. Two roots are equal
by (2). Finally $\pm2e_i-r$ cannot belong to $s\Z^n$ for any
$s\ge2$: at least one coordinate belonging to the other coordinate
of $r$ has value $\pm1$. Thus the sum is the same root, and the
two possible intermediate vertices are exactly the other square
corners. Translation of the entire argument gives the dual statement.
\end{proof}

We will also use the following form of local lifting. Fix a unit cube
$C$, or a cube in the dual grid. Its vertices and subfaces inject into
the quotient: a nonzero period cannot move a vertex or a subface of
$C$ to another vertex or subface of $C$, since the coordinate
differences have absolute value at most one. An incident chain of
face orbits has a common lift by first fixing its largest face and
then lifting the smaller faces inside it. These lifts are unique.
Incidences between subfaces of that largest face are consequently
literal containments. When a vertex and a chain of faces form a flag,
the vertex has its unique lift inside the smallest face. These
observations apply to primal and dual cubes, and will be used together
with Lemma~\ref{NI.separation} when a flag has two distinct vertex types.

\begin{proposition}\label{NI.first}
The original left and right leaves satisfy B1 and B2 exactly when
$(k,s)\ne(1,2)$. After the first left operation the geometry $H_{10}$
has the following literal model modulo $L$:
\[
\begin{array}{c|l}
0&\text{even integral vertices}\\
1&\text{odd integral vertices}\\
i\ge2&\text{original }i\text{-cubes}.
\end{array}
\]
Incidence between types 0 and 1 is unit-step adjacency; all incidences
with cubes and between cubes are containment. This geometry is a
regular hypertope.
\end{proposition}
\begin{proof}
Unit-step edges can become parallel in the quotient only if two
distinct signed unit vectors differ by a period. Because periods are
in $s\Z^n$, this is exactly the case $2e_i\in\Lambda_k$, settled
by Lemma~\ref{NI.separation}(1). To prove B2, lift a cube and two of its
vertices. If their difference is a unit vector modulo $\Lambda_k$,
the error has coordinates of absolute value at most two and can have
a nonzero even coordinate only on that single axis. The only possible
period is again $\pm2e_i$ at $(k,s)=(1,2)$. Otherwise the edge is a
literal edge of the cube. The forward implication is containment.
Self-duality proves the right-end assertion.

The parity of the coordinate sum colors the infinite grid. Its
color-preserving period subgroup is exactly $L$. The component
construction on the quotient by $\Lambda_k$ is consequently the
displayed colored-vertex construction modulo $L$. More explicitly,
if $\Lambda_k\subset D$ the quotient grid is bipartite and no
objects are doubled; if an odd period exists, there are two cosets
of $L$ in $\Lambda_k$, giving the two vertex lifts and the two
compatible lifts of each cube. The three displayed lattices lie in
$D$ exactly when $ks$ is even; an odd period otherwise supplies
an odd closed walk in the quotient grid. Every cube shadow is an ordinary
cube graph and is bipartite; higher flags are chains of cubes and
their shadows are the smallest cubes. Thus PC holds, including
in the globally nonbipartite case. The retained halving theorem
gives thinness, residual connectivity and regularity.
\end{proof}

\section{A second operation at one end}
\begin{proposition}\label{NI.second}
In $H_{10}$ the leaves $(0,2)$ and $(1,2)$ satisfy B1 and B2
exactly when $(k,s)\ne(2,2)$, assuming $(k,s)\ne(1,2)$ already.
Under this condition the second-left graph is nonbipartite, PC
holds, and its halving $H_{20}$ is a regular hypertope. Its objects
modulo $L$ are
\[
\begin{array}{c|l}
0,2&\text{two separate copies of the even integral vertices}\\
1&\text{odd integral vertices}\\
i\ge3&\text{original }i\text{-cubes}.
\end{array}
\]
The two even-vertex types are incident when their underlying vertices
differ by a root in $R$; even and odd types are incident at unit
distance; incidences with cubes and between cubes are containment.
\end{proposition}
\begin{proof}
Each square has exactly two even corners. Its resulting edge has
a root as difference vector. Conversely a root specifies a unique
square. Lemma~\ref{NI.separation}(2) therefore proves simplicity except
at the stated small period. For B2 at a higher cube use part (3) of
that lemma. At an odd vertex $w$, the even neighbors are the $2n$
points $w\pm e_i$. They are distinct by part (1). Two such neighbors
give a square through $w$ exactly when they are on different axes.
Two opposite neighbors on the same axis cannot differ by a root
modulo $L$, by the last argument in part (3). This proves B2.

At $(k,s)=(2,2)$ the squares on coordinates 1,2 with respective
corner sets $\{0,e_1,e_2,e_1+e_2\}$ and
$\{0,-e_1,-e_2,-e_1-e_2\}$ are distinct, but have the same
two even corners modulo $2D$. Hence B1 fails. B2 fails too:
the odd vertex $e_1$ is adjacent to both these even corners but
is not a corner of the second square modulo $2D$.

The root graph is connected because $R$ spans $D$, and contains
the triangle $0,e_1+e_2,e_1+e_3$. Its three vertices remain
distinct. The graph shadow of an $m$-cube, $m\ge3$, is its
demicube: one parity class of its corners, adjacent at Hamming
distance two. It is connected and contains a triangle. Indeed,
any even difference can be removed two coordinates at a time,
and the three corners with difference supports
$\varnothing,\{1,2\},\{1,3\}$ form a triangle (with a translated
choice for the other parity). The shadow at an odd vertex is the
cross-polytope graph on $\{w\pm e_i\}$, with edges between
nonopposite pairs. It is connected and contains a triangle for
$n\ge3$.

A higher flag is either a chain of cubes, whose common shadow is
the demicube of its smallest member, or such a chain together
with an odd vertex $w$. In the latter case its even neighbors
inside the smallest $m$-cube are the $m$ corners obtained from
$w$ by a single allowed step; they form $K_m$, with $m\ge3$.
If there is no cube the shadow is the full cross-polytope graph.
Local lifting justifies these descriptions in the quotient. All
of them are nonbipartite, proving PC. Thus every higher object has
just one lift in the new component construction, while the even
vertices acquire two copies. This gives precisely the asserted
model and proves its regular-hypertope conclusion.
\end{proof}

There is a useful consequence of this proof. Any flag of just the
three lower types in $H_{20}$ has a common lift whose vertices span
an original face of dimension at most two. A single vertex spans a
point; an even--odd pair spans an edge; a pair of even types spans
their unique square. If the third, odd, type occurs, Lemma~\ref{NI.separation}
puts it at one of the other corners of that square. The span is
unique. If a cube is also prescribed, the span lies inside its
unique lift. For $H_{10}$ the corresponding lower flags span at
most an edge. We call this the \emph{small-face property} below.

\section{Operations at the opposite end}
Cube centers form the dual grid $\frac12\one+\Z^n$. Color a
center $c=\frac12\one+v$ by the parity of $\sum v_i$.
This is well-defined modulo $L$. For a primal $d$-face $F$, the
centers of the $n$-cubes containing $F$ form the vertices of a
dual cube of dimension $n-d$. A facet between two such cubes
is a dual edge. These are literal local statements, also after
quotienting by $L$, by the local-lifting observation.

\begin{proposition}\label{NI.oppositefirst}
One may perform the first right operation after $H_{10}$ for all
$n\ge3$, and after $H_{20}$ for $n\ge4$. The selected graph is
bipartite, B1 and B2 hold, and the resulting geometries $H_{11}$
and $H_{21}$ are regular hypertopes.
Their two upper types $n,n-1$ are respectively the two colors of
original $n$-cube centers. These types are incident precisely when
the cubes share a facet. Lower objects retain their earlier model,
and are incident with an upper cube when they are contained in it.
The cubes previously of dimension $n-1$ are removed.
\end{proposition}
\begin{proof}
In both cases the original $n$-cubes and $(n-1)$-cubes are still
present: for $H_{20}$ this uses $n-1\ge3$. Their graph is the
dual unit-step grid modulo $L$. It is simple by
Lemma~\ref{NI.separation}(1), and bipartite by the displayed coloring.
B2 for a lower original cube is the primal/dual version of the
edge argument in Proposition~\ref{NI.first}. For a lower vertex,
including either copy of an even vertex, it is the same argument
for the dual $n$-cube around that vertex. No other object types
occur. Every higher shadow is a connected bipartite cube graph:
for a flag without an original cube, use its small-face span;
when there is a chain of original cubes, use the largest one.
Their codimension is at least two. The component construction
therefore splits the center graph into its two colors, does not
duplicate any other object, and has precisely the model asserted.
Residual connectivity follows from the bipartite case of the
retained theorem, without an additional PC assumption.
\end{proof}

\begin{proposition}\label{NI.oppositesecond}
A second right operation is possible in $H_{11}$ when $n\ge4$
and in $H_{21}$ when $n\ge5$, provided $(k,s)\ne(2,2)$.
Its graph is nonbipartite, B1 and B2 hold, and PC holds. It
therefore gives $H_{12}$ (dual to $H_{21}$) and $H_{22}$,
respectively, as regular hypertopes.
The two new upper types $n,n-2$ are separate copies of the even
centers, incident at root distance; type $n-1$ consists of odd
centers. All remaining objects keep their earlier meanings.
\end{proposition}
\begin{proof}
The objects of type $n-2$ are still original $(n-2)$-faces under
the rank assumptions. They have four incident $n$-cubes, two
with even centers. Thus the selected graph is the root graph
on the even dual vertices. The proofs of simplicity and B2
are exactly Lemma~\ref{NI.separation} in the dual grid, applied
separately to each actual higher object: a remaining original
face has a dual cube as its shadow, a lower vertex has a dual
$n$-cube, and an odd center has its cross-polytope neighbor graph.
Different copies of a lower vertex have the same singleton shadow,
which causes no problem for B2.

For completeness consider an arbitrary higher flag, not just its
singletons. First suppose it contains no odd center. If it contains
original faces, their largest member $F$ determines the common
dual cube, since all the lower vertices are in $F$. Its dimension
is $n-\dim F\ge3$, because type $n-2$ is being halved and
the remaining face dimensions are at most $n-3$.
If the flag contains only lower vertex types, their small-face
span has dimension $d\le1$ in $H_{11}$ and $d\le2$ in $H_{21}$.
The common even-center shadow is the demicube of a dual cube
of dimension $n-d\ge3$. This also covers a singleton. An empty
flag has the whole root graph, which is connected and nonbipartite.

Now suppose an odd center $c$ is in the flag. Lift its primal
$n$-cube. Every prescribed lower vertex and every prescribed
primal face lifts inside this cube. The small-face property and
local lifting make these lifts mutually compatible. Let $F$ be
their largest face, or their small-face span if there is no
prescribed face. A neighboring even-center cube contains $F$
exactly when its common facet with the cube at $c$ contains $F$.
There are $n-\dim F$ such facets, so the shadow is the clique
$K_{n-\dim F}$. Its size is at least three by the same bounds.
If there are no other objects, the shadow is the full
cross-polytope graph. All these graphs are connected and
nonbipartite. This proves PC for every flag. The component
construction now gives exactly the stated duplicate even-center
model. The retained theorem proves the conclusion.
\end{proof}

These arguments also identify the order of the operations. Where
both are defined, doing the first right operation before the second
left operation yields the same $H_{21}$ incidence system. One can
either apply Proposition~\ref{NI.oppositesecond} by duality to $H_{11}$,
or compare the displayed models: both have two even-vertex copies,
odd vertices, the remaining primal cubes, and the two center colors,
with exactly the incidences already specified. The same comparison
identifies the two paths to $H_{22}$. Thus an arrow in a diagram of
constructions does not rely merely on equality of the group orders.

\section{Exact diagrams and exhaustion of leaf operations}
The regular-marking theorem in~\cite{Coset} replaces a selected leaf
generator $r$ by $rqr$, where $q$ is its neighbor. Therefore the
markings of the models just proved have the following uniform form:
\[
 (\ell_a,s_1,\ldots,s_{n-1},u_b),
\]
where
\[
\begin{array}{c|c|c}
j&\ell_j&u_j\\\hline
0&\tau_{e_1}z_1&z_n\\
1&\tau_{e_1+e_2}z_1z_2(1\,2)&z_{n-1}z_n(n-1\,n)\\
2&\tau_{e_1+e_3}z_1z_3(1\,3)&z_{n-2}z_n(n-2\,n).
\end{array}
\]
Here $(i\,j)$ denotes a coordinate swap. The second row is the
first conjugation at each end; conjugating the next inward swap
gives the third. This also checks all overlap cases in small rank.
Since the geometries are thin and residually connected, these
product orders are the actual rank-two residue gonalities.

We give indexed edge tables to make the small-rank diagrams
unambiguous. Unlisted pairs have order two. Let
\[
 P=\{\{i,i+1\}:1\le i\le n-2\}.
\]
Its edges all have label three. In the next table, an edge listed
in the last column has label four; all other listed edges have
label three. Repetitions in a union do not add another edge.
\[
\begin{array}{c|l|l}
\text{state}&\text{edges of label three}&\text{edges of label four}\\\hline
00&P&01,\ (n-1)n\\
10&P\cup\{02\}&(n-1)n\\
11,\ n=3&P\cup\{02,03,13\}&\varnothing\\
11,\ n\ge4&P\cup\{02,(n-2)n\}&\varnothing\\
20,\ n=3&P\cup\{01,02\}&03,23\\
20,\ n\ge4&P\cup\{01,02,03\}&(n-1)n\\
21,\ n=4&P\cup\{01,02,03,04,24\}&\varnothing\\
21,\ n\ge5&P\cup\{01,02,03,(n-2)n\}&\varnothing\\
22,\ n=5&P\cup\{01,02,03,05,25,35,45\}&\varnothing\\
22,\ n\ge6&P\cup\{01,02,03,(n-3)n,(n-2)n,(n-1)n\}
 &\varnothing
\end{array}
\]
For example, $02$ means the unordered pair $\{0,2\}$.
The states $21$ and $22$ only occur at the ranks in the table.

Here is a direct verification of the product orders, independent of
any drawn diagram. Interior swaps on disjoint supports commute;
consecutive ones have product order three. A signed swap on
coordinates $i,j$ has product order three with a swap sharing
exactly one coordinate, and commutes with a disjoint swap. The
two opposite signed swaps on the same pair commute. A sign change
on one coordinate of a signed swap has product order four.
The same statements apply to the affine left reflections after
translating their common fixed intersection to the origin. In the
few overlapping left--right pairs one can check the displayed
maps directly: for $n=3$, $\ell_1u_1$ has order three and
$\ell_2z_3$ has order four;
for $n=4$, $\ell_2$ and $u_1$ have order three; for $n=5$,
$\ell_2$ and $u_2$ have order three. The linear parts already
have these orders. Thus the orders cannot decrease after a
translation quotient. Pairs declared commuting are commuting
as affine maps, so the table holds for every allowed period.

\begin{theorem}\label{NI.complete}
Exclude $(k,s)=(1,2)$. Up to duality and exchange of the two
equivalent fork leaves, the complete sequence of admissible leaf
halvings has states
\[
 00\longrightarrow10\longrightarrow11,
\]
and, if also $(k,s)\ne(2,2)$,
\[
 10\longrightarrow20,\qquad
 20\longrightarrow21\longleftarrow11\quad(n\ge4),\qquad
 21\longrightarrow22\quad(n\ge5).
\]
The initial arrow uses the bipartition construction if $ks$ is even
and the partition construction if $ks$ is odd. The arrows at the
first opposite end are bipartition operations. Each second operation
at an end is a nonbipartite partition operation.
For $n=3$, states $20$ and $11$ are terminal; for $n=4$, state
$21$ is terminal; for $n\ge5$, state $22$ is terminal.
Every permitted state is a regular hypertope, with the exact diagram
in the table.
\end{theorem}
\begin{proof}
Existence, nondegeneracy, parity and PC for all arrows have been
proved in Propositions~\ref{NI.first}--\ref{NI.oppositesecond} and their
dual forms. The original left endpoint reflection $a$ exchanges
vertex parity and the two left fork markings. At the other end
$z_n$ preserves vertex parity and exchanges center parity; it
conjugates the two right fork markings. Whenever a second operation
at the opposite end is already present, the stated rank bounds
make these endpoint reflections commute with its altered marker.
Thus these equivalences still apply at that stage.
The explicit models identify both orders of independent operations.

It remains to justify exhaustion. The table shows that $00$ has
only its two endpoints as leaves. State $10$ has its two left
fork leaves and its right endpoint. At $(k,s)=(2,2)$ the former
fail B1 by Proposition~\ref{NI.second}; the latter still works.
State $20$ has no leaf when $n=3$, and otherwise only the right
endpoint. State $11$ has no leaf when $n=3$; otherwise its four
fork leaves lead, up to the symmetries already described, to $21$.
State $21$ has no leaf when $n=4$, and otherwise just the two
right fork leaves. Finally every vertex in state $22$ has degree
at least two. These are exactly the states and arrows stated.
At the excluded second-step period, the same dual root coincidence
prevents the fork operations in $11$ as well. There is therefore
no additional leaf operation omitted from the classification.
\end{proof}

This is the precise interpretation of the families in Figures~5--8:
the second-step arrows require the extra small-period exclusion;
the last two graphs are the overlapping-end cases in ranks four
and five, not disjoint copies of the large-rank picture. The table
specifies every bond even where a reproduced figure is difficult
to read. The proof supplies the nondegeneracy and PC checks which
the final paragraph of Section~5.3 leaves implicit. It does not
restore the insufficient exact presentations in Propositions
5.8--5.13: the corrected marked groups and all missing period
relations are treated separately in~\cite{Affine}.

\section{Scope and attribution}
The toroids, the halving operation, the proposed diagram families,
and their attribution to earlier polytope and diagram geometry are
retained from the source. The new argument here is an explicit
coordinate verification of every allowed geometric branch, using
the previously reviewed parity-shadow criterion rather than the
source's overbroad residual-connectivity claim. It is a Section~5
recovery component. It makes no whole-original admission, numerical
compression claim, new conjecture claim or formalization claim.

\section{Flag transitivity does not descend}
Theorem A in the introduction of Piedade and Tranchida,
\emph{Constructing new geometries: a generalized approach to
halving for hypertopes}, arXiv:2405.19050v1, asserts that the halved
geometry is flag-transitive if and only if the original geometry is.
We give a counterexample to the converse. The forward assertion in
the detailed Proposition 3.14 is unaffected. This example is not thin;
it is separate from the regular counterexample to residual connectivity.

\begin{lemma}\label{NT.twist}
Let $R$ be a simple graph on $V$, and let $g$ be an automorphism with
$g^2=1$. Suppose $x$ is not adjacent to $gx$ for every $x\in V$.
Define $Q$ by $x\sim_Q y$ if and only if $x\sim_R gy$.
Then $Q$ is simple and undirected. The canonical bipartite double
covers of $Q$ and $R$ are isomorphic by the identity on the first
shore and by $g$ on the second shore.
\end{lemma}
\begin{proof}
Symmetry follows from $x\sim_R gy$ if and only if $gx\sim_R y$,
and the hypothesis excludes a loop. In the double cover, $x$ on the
first shore is incident with $y$ on the second precisely when
$x\sim_R gy$. Thus the stated shore relabeling is the asserted
type-preserving isomorphism.
\end{proof}

Take $V$ to be the two-element subsets of $\{1,2,3,4,5\}$, let
$R$ have adjacency given by disjointness, and put $g=(12)$.
Every $X\in V$ meets $gX$: if $X$ contains exactly one of $1,2$,
its other element is fixed by $g$, and otherwise $X=gX$.
Lemma~\ref{NT.twist} therefore applies. It is useful to write
\[
 a=\{1,2\},\quad p_i=\{1,i\},\quad q_i=\{2,i\},\quad
 b_i=\{3,4,5\}\setminus\{i\}\qquad(i=3,4,5).
\]
The complete edge list of $Q$ is
\[
 \{p_i,p_j\},\ \{q_i,q_j\}\ (i<j),\qquad
 \{a,b_i\},\ \{b_i,p_i\},\ \{b_i,q_i\}\ (i=3,4,5).
\]
Indeed these are exactly the pairs $X,Y$ for which $X\cap gY$ is
empty, by separating the four displayed kinds of vertices. Thus
$Q$ is a connected cubic graph on ten vertices, with fifteen edges.
The $p_i$ and the $q_i$ each form a triangle; $a$ belongs to no
triangle, since its three neighbors are the pairwise nonadjacent
$b_i$. In particular $Q$ is nonbipartite and not vertex-transitive.

Define a rank-three incidence geometry $\Gamma$ as follows. Its
objects of type $0$ are the vertices of $Q$; its objects of type $1$
are the edges of $Q$, with the usual endpoint incidence. Its objects
of type $2$ are two distinct elements $z_+,z_-$, each incident with
every object of the other two types. Distinct objects of the same
type are not incident.

\begin{proposition}\label{NT.example}
The geometry $\Gamma$ is residually connected and satisfies the
source conditions B1 and B2. Its halving $\mathcal P(\Gamma)$ is
residually connected and flag-transitive, whereas $\Gamma$ is not
flag-transitive. It has $60$ chambers, as does its halving.
\end{proposition}
\begin{proof}
Every flag extends to a chamber $(v,e,z_\epsilon)$ with $v\in e$.
At $z_\epsilon$ the rank-two residue is the vertex-edge incidence
graph of $Q$, hence is connected. At a vertex it is $K_{3,2}$;
at an edge it is $K_{2,2}$. The whole incidence graph is connected
because either type-$2$ object is incident with all lower objects.
These are every residue of rank at least two, proving residual
connectivity. Simplicity of $Q$ is B1. Every type-$2$ object is
incident with both endpoints of every edge and with that edge, so
B2 holds. Flag transitivity would imply vertex transitivity of $Q$,
which its triangle distinction excludes. Finally there are
$15\cdot2\cdot2=60$ chambers.

We now identify the halving directly from its definition. Its types
$0$ and $1$ are two copies of $V$, with $X$ on the first shore
incident with $Y$ on the second when $X\cap gY=\varnothing$.
The shadow of either $z_\epsilon$ is the whole connected
nonbipartite graph $Q$. Its canonical double cover is connected,
so $z_\epsilon$ has one lifted object, incident with every object
of the first two types. Thus $\mathcal P(\Gamma)$ is the join of
the bipartite double cover of $Q$ with $\{z_+,z_-\}$.

For completeness, connectedness of that double cover follows without
an enumeration. In a connected nonbipartite graph, a walk between
two prescribed vertices can have either parity: traverse a path to
an odd cycle, go around that cycle, and return along the same path
to reverse the parity. Lifting these walks connects either choice
of sheet over either vertex. The residues in the new geometry are
therefore connected: at $z_\epsilon$ use this double cover, and at
either lower object the residue is $K_{3,2}$. These again exhaust
the rank-two residues.

After the second-shore relabeling $Y\mapsto gY$, the lower incidence
relation is simply disjointness of two-element subsets. The group
$S_5$ is transitive on ordered pairs of disjoint two-element sets:
map the first pair, the second pair, and their unique remaining
element separately. Consequently its action in the original labels
is
\[
 X_0\mapsto(\sigma X)_0,\qquad
 Y_1\mapsto(g\sigma gY)_1\qquad(\sigma\in S_5).
\]
It is transitive on the thirty incident lower pairs. Independently
interchanging $z_+$ and $z_-$ proves transitivity on all sixty
chambers. Since every flag extends to a chamber, this is flag
transitivity. The chamber count is $30\cdot2=60$.
\end{proof}

All nonempty higher shadows here equal $Q$ and are nonbipartite,
so the corrected parity condition is satisfied. The counterexample
therefore does not exploit the separate failure of unconditional
residual connectivity. Its only target is the reverse implication
in the introductory flag-transitivity equivalence. The converse
fails because an automorphism of the bipartite incidence geometry
may act differently on its two shores, and need not descend to an
automorphism of $Q$. The construction and proof use only elementary
graph and incidence arguments; no novelty or publication-priority
claim is made.

\begin{thebibliography}{9}
\bibitem{PT} C. A. Piedade and P. Tranchida, \emph{Constructing new geometries: a generalized approach to halving for hypertopes}, arXiv:2405.19050v1 (29 May 2024). Only this full text is in the present source contract.
\bibitem{MS} P. McMullen and E. Schulte, \emph{Abstract Regular Polytopes}, Cambridge University Press, 2002, Theorem 6D4 and the standard affine Coxeter realization of cubic tessellations.
\bibitem{OldCommon} Human Compression Project, \emph{Changes of chamber generators and binary translation groups}, accepted Topic C common paper, v2. The whole paper, including its binary-module sections, is retained.
\bibitem{Common} Human Compression Project, \emph{Residues and stabilizers in geometric halving}, accompanying common supplement.
\bibitem{Parity} Human Compression Project, \emph{Geometric halving through graph covers and parity-compatible shadows}, phase27 accepted proof, 10 October 2026; its complete geometric proof is retained in the old common paper and the present residue supplement.
\bibitem{Coset} Human Compression Project, \emph{Stabilizers in geometric halving}, accepted proof reproduced in the accompanying common supplement.
\bibitem{Affine} Human Compression Project, \emph{Affine groups after one and two halvings}, accepted proof reproduced in the affine sections of this application.
\bibitem{Phase30} Human Compression Project, \emph{An actual regular leaf hypertope whose halving loses residual connectivity}, accepted phase30 proof and finite certificate; all seven sections are reworked in this application.
\end{thebibliography}

\end{document}
